\documentclass[12pt]{article}
% Préambule commun à tous les documents extraits de « La Revue CAPES »
\usepackage[T1]{fontenc}
\usepackage[utf8]{inputenc}
\usepackage{lmodern}
\usepackage{amsmath,amssymb,amsthm,amsfonts,mathrsfs}
\usepackage{setspace}
\usepackage{listings}
\usepackage{pgfplots}
\pgfplotsset{compat=1.18}
\usepackage{longtable}
\usepackage{adjustbox}
\usepackage{lettrine}
\usepackage{tkz-tab}
\usepackage{changepage}
\usepackage{pdflscape}
\usepackage{graphicx}
\usepackage{tikz}
\usepackage{xcolor}
\usepackage[a4paper,top=2cm,bottom=2.5cm,left=2.5cm,right=2cm]{geometry}
\usepackage{fancyhdr}
\usepackage{draftwatermark}
\SetWatermarkText{La RevueCapes}
\SetWatermarkLightness{0.9}
\SetWatermarkScale{2}
\usepackage{hyperref}
\hypersetup{colorlinks=true,linkcolor=blue!50!black,urlcolor=blue!50!black}

\definecolor{revue}{RGB}{20,60,120}

\pagestyle{fancy}
\fancyhf{}
\renewcommand{\headrulewidth}{0.4pt}
\setlength{\headheight}{15pt}

% \entete{Matière}{Titre}{Type de document}
\newcommand{\entete}[3]{%
  \fancyhead[L]{\small\textsc{La Revue CAPES} -- #1}%
  \fancyhead[R]{\small #2}%
  \fancyfoot[C]{\thepage}%
  \thispagestyle{plain}%
  \begin{center}
    {\color{revue}\rule{\textwidth}{1.2pt}}\\[4pt]
    {\small\textsc{Concours directs CAP-CEG / CAPES -- Burkina Faso}}\\[6pt]
    {\LARGE\bfseries\color{revue} #2}\\[6pt]
    {\large\itshape #1 -- #3}\\[2pt]
    {\color{revue}\rule{\textwidth}{1.2pt}}
  \end{center}
  \vspace{0.5cm}%
}

% Encadré pour les remarques ajoutées dans les corrigés
\newcommand{\remarque}[1]{\par\medskip\noindent\fcolorbox{revue}{revue!6}{\parbox{\dimexpr\linewidth-2\fboxsep-2\fboxrule}{\textbf{\color{revue}Remarque.} #1}}\par\medskip}


\begin{document}
\entete{Culture générale}{Corrigé -- Session 2021 (mesure nouvelle)}{Correction}
\begin{spacing}{1.5}

\noindent\textbf{Sujet.} « \emph{Pour l'accès et le maintien des filles à l'école, il n'appartient pas seulement à ces dernières de s'adapter aux exigences de la vie scolaire, mais aussi à l'institution scolaire de prendre en charge leurs besoins spécifiques.} » Sous la forme d'une dissertation, développez dans une première partie trois arguments qui traduisent les obstacles à la scolarisation des filles du Burkina Faso, et proposez dans une seconde partie trois (03) solutions à même d'améliorer leur accès et leur maintien scolaires.

\remarque{Ce sujet ne figurait pas parmi les corrigés du document d'origine ; la proposition ci-dessous a été rédigée pour compléter la série.}

\rubrique{1. Analyse du sujet}

\begin{itemize}
\item \textbf{Type de sujet :} la consigne impose un plan en \textbf{deux parties} : I. trois obstacles à la scolarisation des filles ; II. trois solutions pour améliorer leur accès et leur maintien. Chaque argument doit être expliqué et illustré.
\item \textbf{Idée directrice de la citation :} la réussite scolaire des filles n'est pas seulement une affaire d'adaptation individuelle ; l'école (et plus largement la société) doit tenir compte des besoins particuliers des filles.
\item \textbf{Mots-clés :} « accès » (entrer à l'école), « maintien » (y rester jusqu'au bout du cycle), « besoins spécifiques » (sécurité, hygiène menstruelle, protection contre les violences et le harcèlement, soutien face aux grossesses précoces, accompagnement, modèles féminins).
\item \textbf{Problématique :} quels obstacles empêchent les filles burkinabè d'accéder à l'école et d'y rester, et comment l'institution scolaire et la société peuvent-elles y répondre ?
\end{itemize}

\rubrique{2. Plan détaillé}

\begin{enumerate}
\item[I.] \textbf{Trois obstacles à la scolarisation des filles.} 1. Les pesanteurs socioculturelles (préjugés, mariages précoces, travaux domestiques). 2. Les difficultés économiques des familles. 3. Un environnement scolaire peu adapté aux besoins des filles (éloignement, absence de latrines, violences, insécurité).
\item[II.] \textbf{Trois solutions.} 1. Sensibiliser les communautés et appliquer la loi. 2. Soutenir financièrement les familles et les élèves. 3. Adapter l'école aux besoins spécifiques des filles.
\end{enumerate}

\rubrique{3. Proposition de rédaction}

\partie{Introduction}

L'éducation des filles est reconnue comme l'un des leviers les plus puissants du développement : une fille instruite deviendra une femme plus autonome, une mère mieux informée et une citoyenne plus active. Au Burkina Faso, la scolarisation des filles a beaucoup progressé, en particulier dans le primaire, mais de nombreuses filles abandonnent encore l'école avant la fin du post-primaire et du secondaire. Or, comme le souligne l'assertion proposée, « pour l'accès et le maintien des filles à l'école, il n'appartient pas seulement à ces dernières de s'adapter aux exigences de la vie scolaire, mais aussi à l'institution scolaire de prendre en charge leurs besoins spécifiques ». Quels obstacles freinent la scolarisation des filles burkinabè, et quelles solutions peut-on envisager pour améliorer leur accès et leur maintien à l'école ? Nous présenterons d'abord trois obstacles, avant de proposer trois solutions.

\partie{I. Trois obstacles à la scolarisation des filles au Burkina Faso}

\souspartie{1. Les pesanteurs socioculturelles}
Dans certaines familles, on considère encore que l'éducation des garçons est prioritaire, car la fille est appelée à se marier et à quitter sa famille ; investir dans ses études paraît alors peu rentable. Les mariages précoces ou forcés mettent fin à la scolarité de nombreuses adolescentes, notamment en milieu rural. Les filles sont aussi chargées de nombreux travaux domestiques (corvée d'eau et de bois, cuisine, garde des cadets) qui réduisent leur temps d'étude et favorisent l'échec scolaire.

\souspartie{2. Les difficultés économiques des familles}
La scolarisation a un coût : fournitures, tenues, frais divers, transport, logement lorsque l'établissement est éloigné. Quand les ressources sont limitées, les familles font souvent le choix de scolariser les garçons. Certaines filles sont retirées de l'école pour aider au commerce familial, travailler comme aides ménagères en ville ou sur les sites d'orpaillage.

\souspartie{3. Un environnement scolaire peu adapté}
L'éloignement des établissements, surtout au post-primaire, oblige les filles à parcourir de longues distances ou à vivre loin de leurs parents, dans des conditions qui les exposent aux risques de violence et de grossesse précoce. Beaucoup d'écoles manquent de latrines séparées et de points d'eau, ce qui pousse les adolescentes à s'absenter pendant leurs menstrues. Le harcèlement et les violences sexuelles en milieu scolaire, ainsi que l'insécurité qui a fermé de nombreuses écoles, aggravent encore la situation.

\transition{Ces obstacles montrent que la réussite scolaire des filles ne dépend pas seulement de leurs efforts personnels : elle exige une action de la société et de l'école.}

\partie{II. Trois solutions pour améliorer l'accès et le maintien des filles à l'école}

\souspartie{1. Sensibiliser les communautés et appliquer la loi}
Des campagnes de sensibilisation, menées avec les leaders coutumiers et religieux, les associations de mères éducatrices et les médias (radios communautaires en langues nationales), peuvent faire évoluer les mentalités sur l'importance de l'éducation des filles. La loi interdisant le mariage des enfants doit être effectivement appliquée, et les cas de mariage forcé signalés et sanctionnés. Il faut aussi encourager un partage plus équitable des tâches domestiques entre filles et garçons.

\souspartie{2. Soutenir financièrement les familles et les élèves}
La gratuité effective de l'enseignement de base, l'attribution de bourses et de kits scolaires aux filles, les cantines scolaires (notamment les rations à emporter qui incitent les familles à envoyer leurs filles à l'école) et les transferts monétaires aux familles vulnérables réduisent le coût de la scolarisation. Des vélos peuvent être distribués aux élèves qui habitent loin de l'établissement.

\souspartie{3. Adapter l'école aux besoins spécifiques des filles}
C'est là le cœur de l'assertion : l'institution scolaire doit prendre en charge les besoins des filles. Cela passe par la construction d'établissements plus proches, de latrines séparées et de points d'eau, la mise à disposition de protections hygiéniques, la création de foyers ou d'internats sécurisés pour les filles éloignées de leur famille, la lutte sans complaisance contre le harcèlement et les violences en milieu scolaire, l'éducation à la santé sexuelle et reproductive, l'accompagnement des élèves enceintes pour qu'elles puissent reprendre leurs études, et la présence de femmes enseignantes qui servent de modèles.

\partie{Conclusion}

L'accès et le maintien des filles à l'école au Burkina Faso se heurtent à des obstacles socioculturels, économiques et institutionnels : préjugés et mariages précoces, pauvreté des familles, environnement scolaire inadapté et insécurité. Pour les surmonter, il faut sensibiliser les communautés et faire appliquer la loi, soutenir financièrement les familles et les élèves, et surtout adapter l'école aux besoins spécifiques des filles. Comme l'affirme l'assertion, la réussite des filles est une responsabilité partagée entre les filles elles-mêmes, leurs familles et l'institution scolaire. Scolariser et maintenir les filles à l'école, c'est préparer des femmes capables de contribuer pleinement au développement du pays.

\end{spacing}
\end{document}
